\Article{11}
Man stelle sich von jedem ganzen Punkt $k$ nach $k+\frac12i$ gerade Linien gezogen vor, deren rechte Seite als positiv und deren linke als negativ gilt. Es sei $l$ eine Linie, und $Sl$ bezeichne die Summe aller Schnitte von $l$ mit diesem Liniensystem, wobei diejenigen positiv gezählt werden, bei denen $l$ von der negativen auf die positive Seite übergeht, und die entgegengesetzten Schnitte negativ. Für den Dezidenten gilt dann der folgende Ausdruck:
\[
\Sigma(\Upsilon l\pdot Sl)+\Sigma Sl'-S\mu
\]

Dabei bezeichnet $l$ alle Ligaturen der Quadrate in $\omega$ (stets so genommen, dass die Quadrate ihnen zur Rechten liegen), und $l'$ diejenigen Ligaturen, die auf dem Umfang der Figur $\omega$ zwischen $0$ und $\frac12m$ liegen, also außerhalb von $\Omega'$.

Alle Ligaturen $l$ bestehen aus
\begin{enumerate}
\item[1)] $l'$;
\item[2)] $l''$, den innerhalb von $\Omega'$ liegenden Grenzligaturen, also $\lambda',\overline{\overline\nu},\overline{\overline\lambda}$;
\item[3)] $l'''$, die im Innern von $\omega$ liegen.
\end{enumerate}
Verstünde man unter $l$ unbestimmt alle Ligaturen, die innerhalb von $\omega$ oder auf den Grenzen dieser Figur liegen, sofern sie von Punkten $\frac{km}{2M}$ ausgehen, wobei $k$ durch $1+i$ teilbar ist, so wäre der Dezident
\[
=\Sigma\alpha\pdot Sl-S\mu
\]

\begin{list}{}{\setlength{\labelwidth}{5em}\setlength{\leftmargin}{5.5em}}
\item[wobei $\alpha=1$] für alle Ligaturen im Innern von $\omega$;
\item[$\alpha=\Upsilon l+1$] für alle Grenzligaturen außerhalb von $\Omega'$, deren Richtung entlang der von $0$ nach $\frac12m$ gehenden Grenze verläuft;
\end{list}

\SourceLeaf{341}

\begin{list}{}{\setlength{\labelwidth}{9em}\setlength{\leftmargin}{9.5em}}
\item[$\alpha=-(\Upsilon\overline{\overline\lambda}+1)=-\Upsilon l$] für alle auf dieser Grenze, die in entgegengesetzter Richtung verlaufen;
\item[$\alpha=\Upsilon l$] für alle Grenzligaturen innerhalb von $\Omega'$, deren Richtung auf $0$ zugeht;
\item[$\alpha=-\Upsilon l+1$] für alle Grenzligaturen innerhalb von $\Omega'$, deren Richtung von $0$ wegführt.
\end{list}

\Article{12}
Wir können nun sämtliche vorkommenden $l$ (nach dem letzten Verfahren) paarweise zusammenfassen, nämlich $l$ mit $\overline{\overline{\frac12m-l}}$, die wir verbundene Ligaturen nennen wollen. Eine einzige ist davon ausgenommen, die isoliert steht oder mit ihrer verbundenen Ligatur identisch ist, nämlich diejenige, die von
\[
\tfrac12(M-1)\pdot\frac{m}{2M}\quad\text{nach}\quad\tfrac12(M+1)\frac{m}{2M}\quad\text{läuft}
\]
Für verbundene Ligaturen ist $\alpha$ stets gleich.
\BookRule{.3\textwidth}
